\section{How this paper and its website were made}\label{app:making-of}

This appendix records how the paper and its companion site were produced: who did what, how the material was found, how the results were generalised, unified and checked, and what it cost.

\subsection{In brief}

\begin{itemize}
\item 1 October 2026, 17:06 AEST, to 4 October 2026, about 01:00 AEST, when this appendix was last revised: about 56 hours of wall-clock time. At least one agent was running for about 23 of those hours. The rest was a 43-minute wait for an account usage reset, two interruptions, and long pauses between requests.
\item The authors directed one orchestrating session of Claude (Anthropic's Claude Opus 5.5) running in Claude Code, in short messages and a few answers to multiple-choice questions. That session wrote and launched 20 multi-agent workflows, which did most of the research, writing and checking, and two small ones that checked and redrafted this appendix. It did the planning, integration and publishing itself.
\item Those workflows started 422 agent runs (420 distinct agent calls, two of which were started twice). 358 completed, 29 failed when the account reached a spend limit, and 35 were stopped or cut off.
\item The agents processed about 117 million fresh tokens: prompts written to the cache plus generated text, of which 8.6 million were generated. They also re-read 4.1 billion tokens from the prompt cache. The orchestrating session added about 10 million fresh tokens and 0.2 billion cache reads.
\item The outputs are a website (15 sections including a version of this appendix, 20 interactive figures plus an animated hero, and a checked catalogue of 382 methods and 138 papers) and a paper of about 400 pages with the same numbering.
\end{itemize}

\subsection{Who did what}

The authors think that when the literature on parameter-efficient fine-tuning (PEFT) already holds several hundred methods, the field gains more from the approach Grothendieck pictured as a rising sea than from one more method. This atlas catalogues 382 of them.

In his memoir \emph{Récoltes et Semailles}, written 1983--1986, Grothendieck compared a hard problem to a nut. One can strike the shell with hammer and chisel, at point after point, until it cracks. Or one can soak the nut in a softening liquid and let time pass, until the pressure of a hand is enough. On a larger scale the second way is the rising sea. It advances so slowly and quietly that nothing seems to change, until it surrounds the problem and the problem is dissolved by a theory that reaches well beyond it~\citep{xgrothendieck2022recoltes,xmclarty2007rising}. In the authors' analogy, a 383rd method would be one more strike of the chisel, and the sea is a theory in which the existing methods appear as instances of one notion.

The authors have packaged this way of working as a Claude Code skill. It asks for the natural home of a result, for objects defined by their relationships and for the relative point of view, so that hard results dissolve. Its binding rule is that every reframe must pay, by a shorter proof, a dropped assumption or a real unity, or be rejected.

This view shaped the question the authors chose, to re-found PEFT ``as Grothendieck would'' and to make the result an interactive site. They also made the main decisions of scope and presentation. Their instructions to the session were short messages: pointing out some missing methods in the literature, asking for a sweep of 2026 work, a mathematical appendix, more readable typefaces, a LaTeX paper and this appendix, supplying the author list and the title, and asking the session to resume after interruptions.

The orchestrating session did the following:

\begin{itemize}
\item planned the work and wrote the seed sketch of the theory;
\item designed the visual identity, the data pipeline, the build and test tools and the LaTeX template;
\item launched and resumed the workflows, read their reports and decided what to run next;
\item applied some changes itself, such as the new typography rules, and rebuilt, tested and published the site, and built the paper.
\end{itemize}

The sub-agents researched the literature, proposed, judged and refereed the theory, classified the methods, wrote the sections, drew and tested the figures, and reviewed each other's work. They also made the remaining decisions, including which 2026 candidates entered the catalogue and what was cut from the site in its final review.

A workflow is a script for Claude Code's Workflow tool. It starts sub-agents in parallel or in sequence, hands each a prompt and files, and keeps a journal of their results. Every sub-agent ran the same model, \texttt{claude-opus-5-5}, whose training data ends in June 2026. All had web search and fetch and a shell on one Apple M5 laptop, with Python 3.9 (NumPy, PyTorch, SciPy, Matplotlib), Node.js, headless Chrome and TeX Live 2026. Below, an independent check means one made by a separate agent with its own context and instructions.

\subsection{Finding the material}

The literature work was split into thirteen families:

\begin{itemize}
\item the Hugging Face PEFT library;
\item LoRA and its training-dynamics variants;
\item initialisation and rank allocation;
\item sharing and extreme efficiency;
\item structured shapes;
\item multiplicative and orthogonal methods;
\item adapters;
\item prompts and states;
\item selective fine-tuning;
\item quantisation and memory;
\item composition and merging;
\item theory and foundations;
\item the 2025--2026 frontier.
\end{itemize}

One agent swept each family. It worked from the arXiv API, Semantic Scholar, OpenReview, the ACL Anthology, the official conference listings and the Hugging Face PEFT source code. Each sweeper was told to read a paper's method section before writing down its formula. As each family finished, an independent fact-checker re-verified every entry:

\begin{itemize}
\item title, first author and first-version year against the arXiv API;
\item the venue against the proceedings;
\item the formula against the method section;
\item the parameter count, re-derived independently;
\item mergeability and initialisation;
\item Hugging Face PEFT support, against the repository at commit 85be3f3 (30 September 2026).
\end{itemize}

A completeness critic then compared the whole catalogue with the library and with its own knowledge of the field. It listed 34 missing methods, and 32 of them were researched and checked in four more batches. The checkers, including those of the four gap-filling batches, made 613 corrections before the first release. Merging duplicates on a canonical identifier and an alias table gave 320 methods and 124 papers. Matching by name would have folded LoRA+ into LoRA.

After the first build, the authors pointed out methods missing from the catalogue. They were added and refereed, and a dedicated sweep of 2026 work followed, in which six scouts covered:

\begin{itemize}
\item arXiv submissions from January to September 2026, one scout per quarter, by keyword queries over titles and abstracts;
\item the accepted papers of ICLR 2026 and ICML 2026, with ACL and CVPR 2026;
\item the citation graph of ten anchor PEFT papers and the 2026 Hugging Face PEFT releases.
\end{itemize}

A triage agent then chose which candidates entered the atlas. It admitted new 2026 methods with at least one signal (a 2026 venue, inclusion in Hugging Face PEFT, citations, a widely used codebase, or a structure the seven coordinates would treat differently), at most 36 per pass. The first triage saw only the first 105 of the 265 candidates not already in the atlas, the last of them cut off mid-entry, because the orchestrator had truncated the list it passed in. A second pass triaged the remaining 161 candidates and the scouts' overflow lists. Together the passes added 61 methods and 14 theory papers. They were researched and fact-checked in the same way, giving 382 methods, 138 papers and 905 corrections in total.

The BibTeX of the 502 catalogue entries is built from arXiv metadata for the 493 that have an arXiv identifier, and the other nine are built from the catalogue's own fields. Eight of those nine were corrected by hand for the paper. The paper's 130 other references, mostly mathematics and related work, were written by the agents who cited them, who were told to check every field against the publisher, DOI or arXiv page.

\subsection{From examples to a theory}

The orchestrator wrote a seed sketch of eleven numbered points, about 1,000 words, which the architects could adopt, correct or reject. Its points were:

\begin{itemize}
\item methods as reparametrisations in Para and as objects of a pointed slice;
\item frozen and full fine-tuning as initial and terminal;
\item mergeability as a lifting problem;
\item Klein's Erlangen programme for group-type methods;
\item the $\GL_r$ gauge of LoRA and its conserved charge;
\item string diagrams for structured updates;
\item a monoidal sum;
\item base change;
\item the lens duality between forward and backward methods;
\item prompts and prefixes as learned input states;
\item naturality, so that a method defined once per layer type extends to every network.
\end{itemize}

Four architect agents each developed a complete framework from this sketch, with definitions, propositions and proofs, from a different point of view:

\begin{itemize}
\item parametric lenses;
\item Grothendieck's relative point of view;
\item the Erlangen programme with gauge symmetry;
\item string diagrams.
\end{itemize}

Three judge agents scored all four proposals, each from an assigned point of view:

\begin{itemize}
\item a category theorist's, for rigour;
\item a PEFT researcher's, for faithfulness to the methods;
\item a science communicator's, for clarity and visual potential.
\end{itemize}

All three recommended the relative point of view as the backbone. One synthesis agent then built the framework on it. As the judges advised, it took the lens and cometric from the parametric-lens proposal, the covariance and optimizer gauge from the Erlangen proposal and the diagrams from the string-diagram proposal, and it cut the constructions they judged decorative. The standing rule was that a categorical construction must explain, classify or predict something about PEFT, or be cut. The result was 41 numbered results, a seven-axis classification and a dictionary between the two vocabularies.

Generalisation usually followed one pattern:

\begin{enumerate}
\item Take an observation about a single method. Zero-initialised LoRA cannot move in some directions at the first step.
\item Find the invariant it depends on. Here it is the first-order image compared with the dimension of the image germ.
\item State the result for every method with that coordinate value. Any method whose base point is the vertex of a cone has the same deficit.
\item Look for consequences the source papers do not state. Householder reflection adaptation at its default start is also an apex.
\item Test the claim numerically and send it to referees.
\end{enumerate}

The same pattern turned ``LoRA+ helps'' into an exact statement about Adam's invariance to per-block rescaling. It turned ``orthogonal methods preserve structure'' into complete invariants of each group action. It turned ``this adapter can be merged'' into a local rewriting test. The test certifies exact merging under stated side conditions, names the weight that absorbs the update, and depends on the activation, the normalisation and the bias slots.

The seven coordinates are decidable from a method's formula. That made it possible to classify the first 320 methods in ten batches of 32, and the 62 later additions in three more. In each batch a second agent re-derived every coordinate, checked every arrow's factorisation and base point by hand, and spot-checked a few gauge dimensions and first-order deficits numerically. A consolidator then reconciled the first ten batches with each other and with the theory, and the final site review made the prose agree with the data. Every arrow between methods records its explicit map h with $\rho_N\circ h=\rho_M$ (403 of them). Every recorded non-arrow names the test that rules it out (485).

\subsection{How category theory unifies the field here}

The unifying move is to study each method through its map into weight space. A weight-space method is a pointed map from its small parameter space into weight space that sends its starting point to the pretrained weights. Comparing methods then means comparing maps:

\begin{itemize}
\item a method N simulates a method M when M's map factors through N's, that is $\rho_N\circ h=\rho_M$ for a pointed map h;
\item doing nothing and full fine-tuning are the initial and terminal objects, so they bracket every weight-space method;
\item three invariants of the map carry most of the content: its image germ, its fibres with the optimizer's metric, and its behaviour under base change.
\end{itemize}

Extensions such as adapters and prefixes, which change the network itself, fit through a realisation map. Mergeability becomes the question of whether a lift exists.

Exposé XI lists what is analogy and what is theorem. The words Noether, Erlangen, gauge and Grothendieck are used in a precise, limited sense, and the paper says which.

\subsection{How the results were checked}

Verification ran in layers:

\begin{itemize}
\item \textbf{Sources.} Every catalogue entry was checked by an agent other than the one that wrote it, against arXiv metadata, the paper and the library source.
\item \textbf{Adversarial model referees, three rounds.} Each of the 41 results went to two referee agents with different lenses: mathematical correctness, and faithfulness to the papers and code. They were told to default to ``refuted'' when in doubt.
\begin{itemize}
\item Round 1 produced 82 verdicts, 45 of them major objections (41 refutations). 27 results were repaired and none was false outright.
\item Round 2 re-examined 26 of the 27 repaired results in 48 reports. Six of its 54 runs were stopped by the spend limit, so Corollary III.13 received no second-round report and four results received one; in the final site review the section editors took the missing referees' place. The reports raised 28 major objections against 18 results, and 19 results were narrowed again.
\item Each round-2 report ended with the strongest statement its referee would sign. These became one public statement per result in \texttt{propositions.json}. The writers and referees of the site and the paper were told to state nothing stronger, and checked for it.
\item A third round, four mathematical referees during the writing of the paper, made 63 further fixes.
\item Those fixes were then carried back into the site, the data and the theory files, 108 edits in all.
\end{itemize}
\item \textbf{Numerical checks.} \texttt{framework\_checks.py}, one of the paper's data files, re-runs most claims marked [Num] in about a minute. The paper's Appendix E lists the few that rest only on the referees' own replications. Each interactive figure computes its numbers in the browser. The engineers of the theory-driven figures checked their key numbers against \texttt{framework\_checks.py} or a separate NumPy computation, and in the final site review a second agent spot-checked the key numbers of every figure.
\item \textbf{Engineering checks.} Every figure was rendered headlessly in light and dark themes at desktop and phone widths, and its builder looked at the images. A whole-site test drives Chrome to check console errors, unmounted figures, overflow, MathJax errors and broken in-page links. The paper is compiled from a clean copy of its sources with pdfLaTeX alone and checked for errors and undefined references.
\item \textbf{This appendix.} Two agents checked every number and date in it against the workflow journals, the agents' transcripts and the project files, and its text was corrected where they disagreed. Its opening was then redrafted by three agents writing independently, merged by a fourth and checked by two more, one for facts and one for style.
\end{itemize}

These checks do not amount to human peer review. Every proposer, writer and referee was the same model, so their errors can be correlated, and the referees of the paper also edited the files they reviewed. A numerical check is not a proof, and every check runs on small float64 examples. No fine-tuning run on a real model was made, so the eleven predictions of Exposé X are untested. As Exposé XI records, the dimension equalities for LoHa and, in general, for VeRA are checked numerically; one clause of corollary 4 of Theorem V.3 is labelled a proof sketch, although Appendix C gives an argument; and Conjectures A, C and D are open. The catalogue is a selection and can still miss methods. Its 2026 additions came from keyword queries on arXiv, four conference lists and one citation graph, filtered by a triage agent, and the completeness critic's own knowledge of the field ends with the model's training data. In the catalogue, 51 method entries and 7 paper entries carry a medium-confidence flag. Some methods were missing until the authors noticed them. The authors are responsible for what the site and the paper say.

\subsection{Building the site and the paper}

The site is a static page. Its prose lives in fifteen HTML sections, counting a version of this appendix, and its data is generated from the verified catalogue and the theory files. Each of its 21 figure scripts was written by an engineer agent (four were finished by a second agent after the spend limit cut off the first), self-tested and later reviewed by another. The page runs from any static host, straight from disk and inside a sandboxed artifact. Its data is a script, so nothing is fetched at run time apart from D3, MathJax and the fonts, which come from public CDNs.

The typography was changed once, at the authors' request. The session set one passage of the site in four typeface schemes, and the authors chose the one it recommended, with Source Serif 4 for text, IBM Plex Sans for labels and IBM Plex Mono for numbers. Seven engineer agents then raised the text inside all 21 figure modules to at least 12 pixels where a reader must read it, and 11 pixels for tick labels. Later the authors found the display face of the title unattractive, and the title was reset in Source Serif 4 as well.

The paper uses a template designed for it: SGA-style exposés with Roman numerals, so that every result has the same number on the site and in the paper. Twenty-one writer agents converted the full 35,000-word theory into LaTeX, with complete proofs in Appendix C. Six figure engineers redrew the key figures as vector PDFs in the same typefaces. An assembler, four mathematical referees, a typesetting referee and a finalizer took it to a clean build of 393 pages, about 400 with this appendix.

\subsection{Cost and interruptions}

Table~\ref{tab:making-cost} breaks the agent runs and tokens down by phase.

\begin{table}[htbp]\centering\caption{Agent runs and tokens by phase. Fresh tokens are prompts written to the cache plus generated text; re-reads from the prompt cache are not counted.}\label{tab:making-cost}\footnotesize
\begin{tabular}{@{}p{5.4cm}rrrr@{}}\toprule\headrow \textbf{Phase} & \textbf{Workflows} & \textbf{Agent runs} & \textbf{Fresh tokens} & \textbf{Generated tokens}\\\midrule
Literature catalogue (first sweep, additions, 2026 sweep) & 4 & 75 & 24.6 M & 0.43 M\\
Theory and adversarial review (rounds 1--2) & 2 & 145 & 16.3 M & 2.03 M\\
Classification, site sections and figures & 6 & 103 & 32.1 M & 2.74 M\\
Site review, appendix and typography & 4 & 49 & 22.1 M & 1.68 M\\
Paper (writing, referees, sync) & 4 & 50 & 22.0 M & 1.75 M\\
\textbf{All workflows} & \textbf{20} & \textbf{422} & \textbf{117.1 M} & \textbf{8.63 M}\\\bottomrule\end{tabular}\end{table}

The table counts all 422 started runs, one transcript each. Cache re-reads (4.1 billion tokens) are not counted, and neither are the orchestrating session and the two later workflows that checked and redrafted this appendix (nine agents).

The work was interrupted three times:

\begin{itemize}
\item at 20:52 on 1 October an account spend limit made 29 agent runs fail, 12 of them mid-run and 17 within seconds of starting;
\item at about 20:40 on 2 October the host machine crashed while the paper was being written;
\item at about 23:10 on 2 October the session was cut off again, and it resumed at 02:04 the next morning.
\end{itemize}

Each time, the orchestrator inspected the workflow journals and launched fresh workflows containing only the unfinished items, telling each agent to read and complete its predecessor's partial file. The six stopped second-round referee runs were not repeated; the final site review covered them instead. A first attempt to resume in place re-ran agents that had already finished. The workflow tool replays only the longest unchanged prefix of agent calls, so that attempt was stopped within two minutes. The other process lessons were:

\begin{itemize}
\item pass large inputs to agents through files, never as truncated prompt text;
\item refereeing a repaired result again is worth it;
\item corrections found while writing the paper must flow back to the site.
\end{itemize}

\subsection{Reproducibility}

The catalogue and the checks accompany the paper as data files:

\begin{itemize}
\item \texttt{catalogue.csv} and \texttt{catalogue.json}, with \texttt{arrows.csv} and \texttt{obstructions.csv};
\item \texttt{propositions.json}, which holds each numbered result with its statement, proof, public statement and the notes of both review rounds;
\item \texttt{framework\_checks.py}, the numerical checks.
\end{itemize}

The LaTeX source carries the text and the generated tables, and the site's figure code can be read in the page. The rest of the record stays with the authors:

\begin{itemize}
\item the raw and verified literature files and the classification batches, in \texttt{research/};
\item the full theory text and the round-2 reports, in \texttt{theory/};
\item the build and test scripts, in \texttt{site/build}, \texttt{site/test} and \texttt{paper/tools}, and the scripts of the paper's figures, in \texttt{paper/figures/src};
\item the 20 workflow scripts, which contain every agent prompt, and the agents' transcripts.
\end{itemize}
